% Generated by openai gpt-4.1 on 2026-10-19T22:19:49.596765+00:00, prompt geosmile-hull-v9 + card-mode addendum, temperature 0 (generate_card_narrative.py --stakeholder manager). Checks: passed. Do not edit by hand.
\renewcommand{\cardHeadline}{\kt{Crime history} dominates the model's reasoning; \kt{Old Town} is the main sensitivity exception.}
\renewcommand{\cardSecA}{Panel 02 shows \kt{crime history} leads the model in most of Hull, with census characteristics next. \kt{Bilton Grange} stands out: green space, unemployment, and census features lead, not crime history. Only \pv{24} areas are not crime-led, and these rarely coincide with high geographic sensitivity---the highest ranks \pv{28}. \kt{Old Town} remains the clear exception for location sensitivity.}
\renewcommand{\cardSecB}{\kt{Intersection density}, road density, and nightlife count are among the \kt{set-aside candidates}: each under \pv{1\%}, together \pv{7\%} of local importance. Related count and density features can split weight, so judge them as groups. The most sensitive areas align with \kt{expert-identified regions} (panel 01).}
\renewcommand{\cardFooter}{Findings show what the model uses, not causes or resourcing priorities.}
\renewcommand{\cardModel}{gpt-4.1}
\renewcommand{\cardTileA}{\tile{144/168}{areas led by six-month crime average}}
\renewcommand{\cardTileB}{\tile{2.6x}{Old Town's sensitivity vs next area}}
\renewcommand{\cardTileC}{\tile{10/30}{features the model barely uses}}
\renewcommand{\cardTileD}{\tile{20/20}{top areas stable when neighbours added}}
